\chapter{Workbook Phần 2: 20 Bài Tập Thực Hành Lập Trình Python \& Pandas}

\begin{lythuyetbox}[Mục Tiêu Luyện Tập Phần 2]
Phần này cung cấp \textbf{20 bài tập code Python \& Pandas thực chiến}, bao gồm cả lập trình giải thuật xử lý dữ liệu và thao tác phân tích nâng cao trên các tệp dữ liệu thực tế, giúp bạn làm chủ 100\% các kỹ năng code được yêu cầu trong CV ứng tuyển DA/DS/DE.
\end{lythuyetbox}

\begin{thuchanhbox}[Tệp Dữ Liệu \& Kịch Bản Thực Hành Python/Pandas Trực Tiếp (Bấm Vào Để Mở)]
Toàn bộ các bài tập phân tích trong chương này được liên kết trực tiếp với các tệp dữ liệu và mã nguồn chuẩn mẫu trong dự án:
\begin{itemize}
    \item \href{run:./data/order_items.csv}{\Large\textbf{\color{navyprimary}data/order\_items.csv}} -- 24,837 dòng chi tiết mặt hàng: \texttt{[order\_id, item\_id, category, price, quantity]}.
    \item \href{run:./data/orders.csv}{\Large\textbf{\color{navyprimary}data/orders.csv}} -- 10,000 đơn hàng: \texttt{[order\_id, customer\_id, order\_purchase\_timestamp, order\_status]}.
    \item \href{run:./data/customers.csv}{\Large\textbf{\color{navyprimary}data/customers.csv}} -- 2,000 khách hàng: \texttt{[customer\_id, customer\_name, city, signup\_date]}.
    \item \href{run:./data/input.csv}{\Large\textbf{\color{navyprimary}data/input.csv}} -- 5,017 bản ghi khách hàng chứa giá trị khuyết và ngoại lai để làm sạch.
    \item \href{run:./code/ch03_python/eda_pipeline.py}{\Large\textbf{\color{navyprimary}code/ch03\_python/eda\_pipeline.py}} -- Pipeline EDA hoàn chỉnh phân tích doanh thu MoM, Pareto 80/20.
\end{itemize}
\small\textit{(Xem trực tuyến tại: \url{https://github.com/TontonYuta/DataLearning/tree/main/data} và \url{https://github.com/TontonYuta/DataLearning/tree/main/code/ch03_python})}
\end{thuchanhbox}

\begin{baitapbox}[Hệ Thống 20 Bài Tập Python \& Pandas Thực Chiến]
\begin{enumerate}
    \item \textbf{Bài 1 (Khử Trùng Lặp List)}: Viết hàm Python \texttt{unique\_sorted(lst)} nhận vào một danh sách số nguyên có thứ tự bất kỳ, trả về danh sách các phần tử duy nhất được sắp xếp tăng dần mà không dùng hàm \texttt{set()}.
    \item \textbf{Bài 2 (Đếm Tần Suất Từ Bằng Regex)}: Viết hàm nhận vào một chuỗi văn bản, trả về một dictionary đếm số lần xuất hiện của từng từ (Word Frequency Counter), bỏ qua dấu câu và không phân biệt chữ hoa/thường bằng module \texttt{re}.
    \item \textbf{Bài 3 (Tìm Cặp Số Có Tổng Bằng K -- Two Sum)}: Cho một mảng số nguyên và một số $K$. Viết thuật toán tìm 2 số trong mảng có tổng đúng bằng $K$ với độ phức tạp thời gian tối ưu $O(n)$ sử dụng Hash Map (Dictionary).
    \item \textbf{Bài 4 (Làm Phẳng List Lồng Nhau -- Flatten Nested List)}: Viết hàm đệ quy làm phẳng một danh sách lồng nhau nhiều cấp (ví dụ: \texttt{[1, [2, [3, 4], 5], 6]} thành \texttt{[1, 2, 3, 4, 5, 6]}).
    \item \textbf{Bài 5 (Xử Lý Ngoại Lệ An Toàn)}: Viết hàm \texttt{safe\_divide(a, b)} chuyển đổi cả hai tham số sang float và thực hiện phép chia. Bắt và xử lý riêng biệt các lỗi \texttt{ValueError}, \texttt{ZeroDivisionError}, và \texttt{TypeError}.
    \item \textbf{Bài 6 (Decorator Đo Thời Gian Chạy)}: Viết một decorator \texttt{@benchmark} đo thời gian chạy của hàm sử dụng \lstinline|time.perf_counter()| và in ra cảnh báo nếu thời gian thực thi vượt quá $1.0$ giây.
    \item \textbf{Bài 7 (Generator Sinh Chuỗi Fibonacci)}: Viết một Generator sử dụng từ khóa \texttt{yield} sinh vô hạn dãy số Fibonacci $(0, 1, 1, 2, 3, 5, 8...)$ với mức tiêu thụ RAM luôn cố định $O(1)$.
    \item \textbf{Bài 8 (NumPy Vectorization \& Min-Max Scaling)}: Cho ma trận $A$ kích thước $(1000, 50)$ và vector $B$ kích thước $(50,)$. Hãy dùng NumPy tính tích vô hướng $A \times B$ và chuẩn hóa vector kết quả về đoạn $[0, 1]$ bằng công thức Min-Max Scaling.
    \item \textbf{Bài 9 (Tối Ưu Hóa RAM Trên \texttt{data/order\_items.csv})}: Mở tệp \href{run:./data/order_items.csv}{\texttt{data/order\_items.csv}} (24,837 dòng). Hãy viết hàm tự động ép kiểu cột \texttt{quantity} từ \texttt{int64} xuống kiểu nguyên nhỏ nhất (\texttt{int8} hoặc \texttt{int16}), ép \texttt{price} xuống \texttt{float32}, và chuyển \texttt{category} (5 giá trị duy nhất) sang \texttt{category}. Đo đạc tỷ lệ RAM tiết kiệm được. \textit{Kết quả kiểm tra}: Tiết kiệm khoảng \textbf{34.2\%} bộ nhớ.
    \item \textbf{Bài 10 (Xử Lý Dữ Liệu Khuyết Trên \texttt{data/input.csv})}: Mở tệp \href{run:./data/input.csv}{\texttt{data/input.csv}} (5,017 dòng). Cột \texttt{age} bị khuyết 131 dòng và \texttt{income} khuyết 50 dòng. Viết mã Pandas điền giá trị khuyết của \texttt{income} bằng giá trị trung vị (median) của từng thành phố tương ứng (\texttt{groupby('city')}).
    \item \textbf{Bài 11 (Tách Chuỗi Mã Định Danh Bằng Regex)}: Trong tệp \href{run:./data/order_items.csv}{\texttt{data/order\_items.csv}}, cột \texttt{item\_id} có dạng chuỗi \texttt{"ORD\_000001\_ITM\_1"}. Sử dụng \lstinline|str.extract()| với biểu thức chính quy để tách thành 2 cột mới: \texttt{order\_code} (ví dụ \texttt{"000001"}) và \texttt{item\_seq} (ví dụ \texttt{"1"}).
    \item \textbf{Bài 12 (Time Series Resampling Trên \texttt{data/orders.csv})}: Đọc tệp \href{run:./data/orders.csv}{\texttt{data/orders.csv}} (10,000 dòng). Chuyển đổi cột \texttt{order\_purchase\_timestamp} sang kiểu \texttt{datetime}, đặt làm index và thực hiện gộp nhóm (Resample) theo từng tháng (\texttt{'ME'}), đếm số lượng đơn hàng phát sinh mỗi tháng. \textit{Kết quả kiểm tra}: Tháng 01/2025 có \textbf{672 đơn hàng}, tháng 02/2025 có \textbf{616 đơn hàng}.
    \item \textbf{Bài 13 (Rolling Average 7 Ngày)}: Dựa trên chuỗi số lượng đơn hàng theo ngày từ \texttt{orders.csv}, sử dụng phương thức \lstinline|rolling(window=7).mean()| để tính đường trung bình động 7 ngày nhằm làm mượt các dao động ngắn hạn cuối tuần.
    \item \textbf{Bài 14 (Tách Cột Danh Sách Bằng \texttt{explode()})}: Cho DataFrame có cột \texttt{tags} chứa danh sách: \lstinline|[['Fashion', 'Sale'], ['Electronics']]|. Hãy viết code Pandas dùng hàm \texttt{explode()} để bung danh sách thành các dòng riêng biệt, phục vụ phân tích tần suất thẻ sản phẩm.
    \item \textbf{Bài 15 (Chuyển Đổi Wide \& Long Bằng Pivot/Melt)}: Đọc dữ liệu từ \texttt{order\_items.csv}. Tính tổng doanh thu theo cặp \texttt{[order\_id, category]}. Dùng \texttt{pivot()} để xoay 5 ngành hàng thành 5 cột ngang, sau đó dùng \texttt{melt()} để chuyển ngược lại cấu trúc bảng dọc Tidy Data.
    \item \textbf{Bài 16 (Phát Hiện Giá Trị Ngoại Lai IQR Trên \texttt{data/input.csv})}: Mở tệp \href{run:./data/input.csv}{\texttt{data/input.csv}}. Tính $Q1$, $Q3$ và khoảng biến thiên tứ phân vị $IQR = Q3 - Q1$ của cột thu nhập \texttt{income}. Xác định ngưỡng trên $Q3 + 1.5 \times IQR$ và ngưỡng dưới $Q1 - 1.5 \times IQR$. Lọc ra toàn bộ các bản ghi thu nhập ngoại lai. \textit{Kết quả kiểm tra}: Tìm thấy chính xác \textbf{27 bản ghi ngoại lai}.
    \item \textbf{Bài 17 (Custom Aggregation Trên \texttt{data/order\_items.csv})}: Tính tổng doanh thu (\texttt{GMV = price * quantity}) cho từng danh mục trong 5 ngành hàng (\texttt{Fashion, Sports, Beauty, Home \& Living, Electronics}). Xác định danh mục dẫn đầu về doanh thu. \textit{Kết quả kiểm tra}: Danh mục \textbf{Fashion} đứng đầu với tổng doanh thu \textbf{\$1,291,623.03}.
    \item \textbf{Bài 18 (Khử Trùng Lặp Giữ Bản Ghi Mới Nhất)}: Trong bảng lịch sử trạng thái đơn hàng, viết cú pháp kết hợp \lstinline|sort_values(by=['order_id', 'timestamp'], ascending=[True, False])| và \lstinline|drop_duplicates(subset=['order_id'], keep='first')| để trích xuất trạng thái cuối cùng của từng đơn.
    \item \textbf{Bài 19 (Phân Tích Nối Bảng Với Merge Indicator)}: Thực hiện phép nối ngoài \lstinline|pd.merge(customers, orders, on='customer_id', how='outer', indicator=True)| giữa 2,000 khách hàng trong \href{run:./data/customers.csv}{\texttt{data/customers.csv}} và đơn hàng trong \href{run:./data/orders.csv}{\texttt{data/orders.csv}}. Đếm số lượng khách hàng chưa từng phát sinh đơn hàng nào (\texttt{'left\_only'}). \textit{Kết quả kiểm tra}: Có chính xác \textbf{11 khách hàng} chưa từng mua hàng.
    \item \textbf{Bài 20 (Xử Lý Tệp Dữ Liệu Lớn Bằng Chunksize)}: Viết script Python sử dụng tham số \texttt{chunksize=5000} trong \lstinline|pd.read_csv('data/order_items.csv')| để tính tổng GMV của toàn bộ 24,837 dòng mà lượng RAM sử dụng tại mỗi thời điểm không vượt quá 5MB. \textit{Kết quả kiểm tra}: Tổng doanh thu là \textbf{\$8,275,862.90}.
\end{enumerate}
\end{baitapbox}

\begin{dapanbox}[Đáp Án \& Lời Giải Ngắn Gọn Cho 20 Bài Tập Python]
\begin{itemize}
    \item \textbf{Bài 1}: 
    \begin{lstlisting}[language=Python]
def unique_sorted(lst):
    res = []
    for x in sorted(lst):
        if not res or x != res[-1]: res.append(x)
    return res
    \end{lstlisting}
    \item \textbf{Bài 2}: \lstinline|from collections import Counter; import re|; \lstinline|Counter(re.findall(r'\b\w+\b', text.lower()))|.
    \item \textbf{Bài 3}: 
    \begin{lstlisting}[language=Python]
def two_sum(nums, K):
    seen = {}
    for x in nums:
        if K - x in seen: return (K - x, x)
        seen[x] = True
    return None
    \end{lstlisting}
    \item \textbf{Bài 4}: 
    \begin{lstlisting}[language=Python]
def flatten(lst):
    res = []
    for item in lst:
        res.extend(flatten(item)) if isinstance(item, list) else res.append(item)
    return res
    \end{lstlisting}
    \item \textbf{Bài 5}: Dùng \texttt{try: return float(a) / float(b)} với các block \texttt{except ZeroDivisionError: return "Lỗi chia cho 0"} và \texttt{except (ValueError, TypeError): return "Tham số không hợp lệ"}.
    \item \textbf{Bài 6}: Decorator sử dụng \lstinline|t0 = time.perf_counter()|, chạy hàm, lấy \lstinline|elapsed = time.perf_counter() - t0|. Nếu \lstinline|elapsed > 1.0| thì in cảnh báo tốc độ.
    \item \textbf{Bài 7}: 
    \begin{lstlisting}[language=Python]
def fib_gen():
    a, b = 0, 1
    while True:
        yield a
        a, b = b, a + b
    \end{lstlisting}
    \item \textbf{Bài 8}: \texttt{C = np.dot(A, B)}; \texttt{C\_norm = (C - C.min()) / (C.max() - C.min())}.
    \item \textbf{Bài 9}: Ép kiểu:
    \begin{lstlisting}[language=Python]
df['quantity'] = pd.to_numeric(df['quantity'], downcast='integer')
df['price'] = pd.to_numeric(df['price'], downcast='float')
df['category'] = df['category'].astype('category')
# Dung lượng RAM giảm từ 4.68 MB xuống 3.08 MB (giảm 34.2%)
    \end{lstlisting}
    \item \textbf{Bài 10}: \lstinline|df['income'] = df.groupby('city')['income'].transform(lambda x: x.fillna(x.median()))|. Đảm bảo tính nhất quán thu nhập theo từng khu vực địa lý.
    \item \textbf{Bài 11}: \lstinline|df[['order_code', 'item_seq']] = df['item_id'].str.extract(r'ORD_(\d+)_ITM_(\d+)')|.
    \item \textbf{Bài 12}: \lstinline|orders['dt'] = pd.to_datetime(orders['order_purchase_timestamp'])|; \lstinline|monthly = orders.set_index('dt').resample('ME')['order_id'].count()|. Tháng 01/2025: 672 đơn, tháng 02/2025: 616 đơn.
    \item \textbf{Bài 13}: \lstinline|daily_orders['rolling_7d'] = daily_orders['order_count'].rolling(window=7).mean()|.
    \item \textbf{Bài 14}: \lstinline|df_exploded = df.explode('tags')|. Mỗi phần tử trong list trở thành một dòng dữ liệu riêng.
    \item \textbf{Bài 15}: Dùng \lstinline|piv = df.pivot_table(index='order_id', columns='category', values='price', aggfunc='sum', fill_value=0)|; Sau đó chuyển lại Long format bằng \lstinline|piv.reset_index().melt(id_vars='order_id', var_name='category', value_name='revenue')|.
    \item \textbf{Bài 16}: $Q1 = 39,708.07$, $Q3 = 60,196.47$, $IQR = 20,488.40$. Ngưỡng trên = $90,929.08$, ngưỡng dưới = $8,975.46$. Phát hiện chính xác \textbf{27 bản ghi ngoại lai} thỏa mãn \lstinline|(df['income'] < 8975.46) | (df['income'] > 90929.08)|.
    \item \textbf{Bài 17}: \lstinline|df.groupby('category').apply(lambda x: (x['price'] * x['quantity']).sum())|. Kết quả: Fashion: \$1,291,623.03; Beauty: \$1,279,015.38; Home \& Living: \$1,269,376.92; Sports: \$1,267,088.89; Electronics: \$1,255,866.63.
    \item \textbf{Bài 18}: \lstinline|df.sort_values(by=['order_id', 'timestamp'], ascending=[True, False]).drop_duplicates(subset=['order_id'], keep='first')|.
    \item \textbf{Bài 19}: \lstinline|merged = pd.merge(customers, orders, on='customer_id', how='outer', indicator=True)|; \lstinline|unengaged = merged[merged['_merge'] == 'left_only']['customer_id'].nunique()|. Kết quả: Đúng \textbf{11 khách hàng} chưa từng mua hàng.
    \item \textbf{Bài 20}: 
    \begin{lstlisting}[language=Python]
total_gmv = 0.0
for chunk in pd.read_csv("data/order_items.csv", chunksize=5000):
    total_gmv += (chunk['price'] * chunk['quantity']).sum()
print(f"Tổng GMV: ${total_gmv:,.2f}") # Ket qua: $8,275,862.90
    \end{lstlisting}
\end{itemize}
\end{dapanbox}
